\section{Innovation \& the Division of Labor}

\subsection{Patents \& Copyrights}
Ideas are \textbf{nonrival} with near-zero $MC$ of use $\implies$ efficient \textit{ex post} price is $\approx 0$, but then creators capture nothing \textit{ex ante}.

\textbf{Static--dynamic trade-off:} IP grants a temporary monopoly: DWL today buys innovation incentive tomorrow. Optimal strength balances:
$$\underbrace{\text{monopoly DWL} \times \text{duration}}_{\text{static cost}} \quad \text{vs.} \quad \underbrace{\Delta(\text{innovations}) \times \text{surplus each}}_{\text{dynamic gain}}$$

\textbf{Margins:} length (duration) vs.\ breadth (how close substitutes may come). Broad-short and narrow-long can deliver the same reward with different DWL profiles.

\textbf{Alternatives to IP:} prizes, first-mover advantage, secrecy, complementary sales (concerts vs.\ recordings). Each avoids pricing above $MC$ but has its own incentive distortion.

\textit{TFU trap:} ``perfect copying is efficient'' --- true \textit{ex post} (competitive allocation), false \textit{ex ante} (no creation). See the two-period copying model in the Monopoly entry: as copying $n \to \infty$, creator revenue $\to 0$.

\subsection{Learning by Doing}
Unit cost falls with cumulative output: $c(Q_{\text{cum}})$, $c' < 0$.

\textbf{Within a firm:} forward-looking pricing---produce beyond static optimum since today's output lowers tomorrow's cost (marginal value of output = price $+$ PV of cost reduction). Can rationalize pricing below current $MC$; \textit{not} predation.

\textbf{Across firms (spillovers):} learning leaks to rivals $\implies$ private return $<$ social return $\implies$ genuine externality; underinvestment. Only \textit{with spillovers} is there an infant-industry case, and the subsidy should target the learning margin, not imports.

\subsection{Division of Labor (Smith, Rosen, Becker--Murphy)}
Specialization has increasing returns: fixed cost of acquiring each skill is spread over its utilization.

\textbf{Model sketch:} worker allocates time across tasks; skill investment per task is a fixed cost $\implies$ utilization rate of each skill rises when fewer tasks per worker $\implies$ per-capita output rises with team size $n$ until coordination costs bind:
$$\max_n \; \frac{A(n)}{n} - \text{coordination}(n), \qquad A(n) \text{ output of } n\text{-member team}$$

\textbf{``Division of labor is limited by the extent of the market''} (Smith): larger markets support finer specialization (higher utilization of each specialty). Modern version: cities, trade, and platforms raise returns to specialists.

\textbf{Coordination costs} limit specialization: hold-up, communication, matching specialists. Firms, contracts, and standards economize on these.

\textit{Applications:} physicians specialize more in big cities; growth accelerates as markets integrate (specialization $\implies$ productivity $\implies$ larger market $\implies$ more specialization---increasing returns without any single firm having them).

\textit{TFU trap:} increasing returns from specialization are \textbf{external to the worker's scale, internal to the market's}: no conflict with competitive equilibrium, unlike firm-level IRS.
